\documentclass[10pt,a4paper]{amsart}
\usepackage[margin=19mm]{geometry}
\usepackage{amsmath,amssymb,mathtools}
\usepackage[scr=rsfs,cal=euler]{mathalfa}
\usepackage{xfrac}
\usepackage[unicode,hidelinks,bookmarks=false]{hyperref}
\newcommand{\ie}{i.e.\ }
\newcommand{\cf}{cf.\ }
\newcommand{\resp}{respectively\ }
\newcommand{\Subsection}{\subsection}
\providecommand{\nobreakdash}{\nobreak\hbox{-}}
\setlength{\emergencystretch}{2em}
\multlinegap0pt
\usepackage{bm}
\usepackage{bbm}
\usepackage{dsfont}
\usepackage{xspace}
\usepackage{cancel}
\usepackage{mathtools}
\usepackage{amsthm}
\makeatletter
\@ifundefined{theo}{\newtheorem{theo}{Theo}}{}
\@ifundefined{prop}{\newtheorem{prop}[theo]{Proposition}}{}
\makeatother
\newcommand{\Nodes}{{\mathcal N}}
\newcommand{\N}{{\mathbb N}}
\newcommand{\pool}{\mathrm{pool}}
\renewcommand{\P}{{\mathbb P}}
\title[Stability of local tip pool sizes]{Stability of local tip pool sizes}
\author{Sebastian M\"uller, Isabel Amigo, Alexandre Reiffers-Masson, Santiago Ruano-Rinc\'on}
\date{Source: arXiv:2302.01625v2 (CC BY 4.0)}
\begin{document}
\maketitle

\noindent\emph{Source and licence.} Stability of local tip pool sizes. Version arXiv:2302.01625v2 (CC BY 4.0), distributed under the Creative Commons Attribution 4.0 International licence, as recorded in the arXiv licence metadata for this exact version and shown on the abstract page https://arxiv.org/abs/2302.01625v2. Full rights notice: licenses/arxiv-2302.01625-CC-BY-4.0.txt.\par\smallskip

\noindent\textit{Editorial scope.} Counted unit: the Prop whose source label is \texttt{prop:small-set-poisson} in the article (source0.tex lines 489--512, source label \texttt{prop:small-set-poisson}), delivered together with the article's own complete original proof of it (source0.tex lines 514--574, one \texttt{proof} environment of 61 lines). No mathematics is written, repaired or re-derived: statement and proof are the article's own text line for line. Required context retained uncounted: none. Every definition, display and label of those lines is retained unchanged. Omitted from this package and not redistributed: 1--488, 513--513, 575--1141. Nothing inside a retained excerpt is omitted or altered: every retained body is delivered exactly as the article's lines are written, so no omission receipt of any kind accompanies this package. Citations: the retained excerpts carry no citation command at all, so no bibliography entry of the article is transcribed and no reference is redistributed. Reference adaptation: none was needed and none was made; every reference inside the delivered unit resolves inside it, so no unresolved reference or citation is produced. Printed slips kept literal: no printed word, symbol, hypothesis or number of the counted statement or of its proof was altered. Layout and class: the article's own class is \texttt{amsart} with packages including amsaddr, inputenc, fontenc, amsmath, amsfonts, amssymb, graphicx, url, multirow, caption, subcaption, booktabs; the wrapper is the lane's pinned \texttt{amsart} editorial wrapper at 10pt on A4, which restates the theorem-like environments the retained text needs and adds the article's own macro definitions that the retained text needs (\texttt{\textbackslash N}, \texttt{\textbackslash Nodes}, \texttt{\textbackslash P}, \texttt{\textbackslash pool}), with extra wrapper packages (bm, bbm, dsfont, xspace, cancel, mathtools). The delivered numbering is the unit's own. Assets: no retained span contains a figure environment, a graphics-inclusion command, a PDF-page inclusion, a table or a TikZ picture; no image file is copied into this package, the package declares no asset dependency and no image file is redistributed with it. Licence: version arXiv:2302.01625v2 of this article is distributed under the Creative Commons Attribution 4.0 International licence, as recorded in the arXiv licence metadata for this exact version and shown on the abstract page https://arxiv.org/abs/2302.01625v2. A full rights notice is delivered at licenses/arxiv-2302.01625-CC-BY-4.0.txt. Characters: every retained excerpt is pure ASCII, so the delivered engine, XeLaTeX with its default Latin Modern text font, reports no missing character.
\begin{prop}[Long-gap minorization and small sets]\label{prop:small-set-poisson}
For a fixed choice of $\alpha>0$, let $\mathsf{Y}_n$ denote the augmented embedded state consisting of the observer tip pool, all local tip pools, and the marked configuration of blocks issued in the last delay window:
\[
   \mathsf{Y}_n
   :=
   \Bigl(
      \pool_n^{(o)},
      (\pool_n(i))_{i\in\Nodes},
      \bigl((t_n-t_m,\kappa_m)\bigr)_{m\le n:\, t_m\in(t_n-\Delta,t_n]}
   \Bigr).
\]
Then $(\mathsf{Y}_n)_{n\ge0}$ is a time-homogeneous Markov chain on a standard Borel state space. Moreover, there exists a probability measure $\nu_\ast$ such that for every $M<\infty$, with
\[
   C_M:=\{y:\ V(y)\le M\},
\]
there are integers $m_M\in\N$ and constants $p_M,\widetilde p_M>0$ for which
\[
   P^{m_M+2}(y,\cdot)\ge p_M\,\nu_\ast(\cdot),
   \qquad
   P^{m_M+3}(y,\cdot)\ge \widetilde p_M\,\nu_\ast(\cdot),
   \qquad \forall y\in C_M,
\]
where $P$ is the transition kernel of $(\mathsf{Y}_n)$. In particular, each $C_M$ is a small set and the chain is $\nu_\ast$-irreducible and aperiodic.
\end{prop}

\begin{proof}
Only blocks issued during the most recent interval of length $\Delta$ can still arrive at different nodes in the future. The current local pools record all tips that may still be selected by some node, while the recent marked configuration determines which not-yet-globally-visible blocks and references arrive before the next issuance. Thus the enlarged state determines all pre-issuance local tip pools at the next decision epoch and therefore determines the law of the next state. By the memoryless property of exponential inter-arrival times and the independence of future marks, $(\mathsf{Y}_n)$ is a time-homogeneous Markov chain.

To make the state space explicit, let $\widehat{\mathsf B}:=\{\star\}\sqcup [0,1]$, where $\star$ denotes the genesis block, and let
\[
   \mathsf K
   :=
   \widehat{\mathsf B}\times
   \Bigl(\bigsqcup_{r=0}^k \widehat{\mathsf B}^{\,r}\Bigr)\times
   \Nodes\times [0,\Delta]^N
\]
be the mark space. Equip $\widehat{\mathsf B}$ with any Borel total order extending $\star<x$ for all $x\in[0,1]$, and for $\ell\ge 1$ write
\[
   \widehat{\mathsf B}^{\,\ell}_{\uparrow}
   :=
   \{(b_1,\ldots,b_\ell)\in \widehat{\mathsf B}^{\,\ell}: b_1<\cdots < b_\ell\}.
\]
We use the convention that $\widehat{\mathsf B}^{\,0}_{\uparrow}$ is a singleton.
Encoding each finite tip pool by the increasing tuple of its elements, the state space of $(\mathsf{Y}_n)$ can be identified with the countable disjoint union
\[
   \mathsf Y
   :=
   \bigsqcup_{\ell_o\ge 1}
   \bigsqcup_{\ell_1,\ldots,\ell_N\ge 0}
   \bigsqcup_{m\ge 0}
   \widehat{\mathsf B}^{\,\ell_o}_{\uparrow}
   \times
   \prod_{i=1}^N \widehat{\mathsf B}^{\,\ell_i}_{\uparrow}
   \times \bigl([0,\Delta]\times \mathsf K\bigr)^m,
\]
which is a standard Borel space because each component is a Borel subset of a Polish space and a countable disjoint union of standard Borel spaces is standard Borel.
We take the actual state space to be the measurable subset of $\mathsf Y$ consisting of admissible configurations generated by the protocol dynamics; this subset is again standard Borel.

Let $\nu_\ast$ be the law of the enlarged post-issuance state obtained when one starts from a synchronised configuration with a single common tip and no recent blocks, and then performs two fresh issuances with inter-arrival gaps larger than $\Delta$. This law includes the observer pool, all local pools, and the remaining recent marked configuration after the second issuance. Fix $M<\infty$ and set $m_M:=\max\{1,\lceil M\rceil\}$. If $\mathsf{Y}_n=y\in C_M$, then necessarily $X_n^{(o)}\le M$. Consider the events
\begin{align*}
   E_M&:=\{t_{n+j}-t_{n+j-1}>\Delta,\ \forall j=1,\dots,m_M+2\},\\
   \widetilde E_M&:=\{t_{n+j}-t_{n+j-1}>\Delta,\ \forall j=1,\dots,m_M+3\}.
\end{align*}
Since the inter-arrival gaps are i.i.d.\ exponential with rate $\lambda$, these events have probabilities
\[
   p_M = \P(E_M)=e^{-\lambda\Delta (m_M+2)}>0,
   \qquad
   \widetilde p_M = \P(\widetilde E_M)=e^{-\lambda\Delta (m_M+3)}>0.
\]
On either event, before each of the next issuances all previously issued blocks are visible at every node, so the recent window has emptied and the local, common, and observer tip pools coincide at each decision epoch. Because $k\ge 2$, whenever the current synchronised observer tip count is larger than $1$, the next synchronised issuance removes at least two distinct tips and creates one new tip, so the observer tip count decreases by at least $1$. Once the observer tip count reaches $1$, further synchronised issuances keep it equal to $1$. As $X_n^{(o)}\le M\le m_M$, after at most $m_M$ such synchronised issuances the chain is in a post-issuance one-tip configuration whose unique tip was created during the long-gap event itself and is therefore independent of the initial state $y$. The next two synchronised issuances, still on $E_M$ or $\widetilde E_M$, then produce an enlarged post-state whose law depends only on the fresh marks generated during those last two issuances; by definition, that law is $\nu_\ast$. Therefore
\[
   P^{m_M+2}(y,\cdot)\ge p_M\,\nu_\ast(\cdot),
   \qquad
   P^{m_M+3}(y,\cdot)\ge \widetilde p_M\,\nu_\ast(\cdot),
   \qquad \forall y\in C_M,
\]
so $C_M$ is small.

For $\nu_\ast$-irreducibility, fix an arbitrary initial state $y$ and write $x:=X^{(o)}(y)<\infty$. On the event of $x+2$ consecutive inter-arrival gaps larger than $\Delta$, the same synchronisation argument shows that after at most $x$ such issuances the observer tip count has reached $1$, and that the final two long-gap issuances produce a state with law $\nu_\ast$. Hence for every measurable set $A$,
\[
   P^{x+2}(y,A)\ge e^{-\lambda\Delta(x+2)}\,\nu_\ast(A).
\]
Therefore the chain is $\nu_\ast$-irreducible.

Finally, for any $M\ge 1+\alpha$, the construction of $\nu_\ast$ gives $\nu_\ast(C_M)=1$. The two minorization relations on such a $C_M$ occur at the coprime times $m_M+2$ and $m_M+3$, so $\gcd(m_M+2,m_M+3)=1$. This is the standard strong aperiodicity criterion for an irreducible chain, hence $(\mathsf{Y}_n)$ is aperiodic.
\end{proof}

\end{document}
